\documentclass[12pt,letterpaper]{report}
\usepackage{graphicx}
\usepackage{scrextend}
\usepackage{vmargin}
\usepackage{graphicx}
\usepackage{multirow}
\usepackage[utf8]{inputenc}
\usepackage[spanish]{babel}
\usepackage{multicol}
\usepackage{enumerate}
\usepackage{float}
\usepackage{amsmath, amsthm, amssymb, amsfonts}
\usepackage[usenames]{color}
\parindent=0mm
\pagestyle{empty}
\definecolor{miorange}{rgb}{0.91, 0.43, 0.0}
\begin{document}
\setmargins{2.5cm}      
{1.5cm}                     
{2cm}  
{24cm}                    
{10pt}                          
{1cm}                          
{0pt}                             
{2cm}
\begin{flushright}
\today
\end{flushright}
\vspace{-1.75cm}
\section*{Tarea Clase 10}
\subsection*{Obtenga las soluciones de las siguientes ecuaciones:}
\begin{enumerate}
\item Escribe de todas las formas posibles la ecuación de la recta que pasa por los puntos A(1,2) y B(-2,5)\\
Calculando la pendiente se tiene que:
\begin{equation*}
    m=\frac{y_2-y_1}{x_2-x_1}=\frac{5-2}{-2-1}= \frac{3}{-3}=-1
\end{equation*}
por lo que nuestra ecuacion de la recta toma la siguiente forma:
\begin{equation*}
    y=-x+b
\end{equation*}
calculando el parametro b, se tiene que:
\begin{align*}
    2&=-(1)+b\\
    b&=3
\end{align*}
por lo que la ecuación de la recta es:
\begin{equation*}
    y=-x+3
\end{equation*}
\item Hallar la pendiente y la ordenada en el origen de la recta $3x+2y-7=0$.\\
Despejando la variable y se tiene que:
\begin{align*}
    2y&=7-3x\\
    y&=-\frac{3}{2}+\frac{7}{2}
\end{align*}
por lo que la pendiente es igual a -3/2
\item La recta $r=3x+ny-7=0$ pasa por el punto A(3, 2) y es paralela a la recta $s=mx+2y-13=0$. Calcula m y n.\\
Sustituyendo los puntos de la recta r se tiene que:
\begin{align*}
    3(3)+n(2)-7&=0\\
    9+2n-7&=0\\
    2n+2&=0\\
    n&=-1
\end{align*}
despejando y la recta s se tiene que:
\begin{align*}
    mx+2y&=13\\
    2y&=13-mx\\
    y&=\frac{13}{2}-\frac{m}{2}x
\end{align*}
como la recta r y s tienen la misma pendiente dado que son paralelas, entonces:
\begin{align*}
    \frac{3}{n}&=-\frac{m}{2}\\
    \frac{3}{-1}&=-\frac{m}{2}\\
    m&=6
\end{align*}
\end{enumerate}
\end{document}